\section{A választott technológiai stack áttekintése}

A megvalósított alkalmazás TypeScript alapú Next.js keretrendszerre épül, amely a React komponensalapú felépítését szerveroldali rendereléssel, fájlrendszer-alapú útvonalkezeléssel és API útvonalakkal egészíti ki \cite{nextjsAppRouter,reactDescribingUi}. Ennek köszönhetően a frontend és a backend jellegű logika egy közös projektstruktúrában jelenik meg: a felhasználói felület React komponensekből épül fel, míg az adatok lekérése és módosítása Next.js API route-okon keresztül történik.

Az alkalmazás adatkezeléséhez és autentikációjához Supabase szolgáltatást használ. A Supabase kliens külön szerveroldali és böngészőoldali inicializálással jelenik meg, ami illeszkedik a Next.js szerveroldali működéséhez és a felhasználói munkamenetek kezeléséhez \cite{supabaseNextjsServerSideAuth}. A rendszer saját adatbázisa elsősorban a felhasználói állapotokat tárolja, például a watchlist és watched listák elemeit, míg a filmek részletes adatai külső forrásból, a TMDB API-ból érkeznek \cite{tmdbApiGettingStarted}.

A kliensoldali adatlekéréseket és a cache-t a TanStack Query kezeli, amely az aszinkron szerverállapot frissítését és szinkronizálását támogatja React alkalmazásokban \cite{tanstackQueryOverview}. A felhasználói felület kialakításához Tailwind CSS, shadcn/Radix UI komponensek és lucide-react ikonok kerültek felhasználásra. A Tailwind utility-first megközelítést biztosít a stílusok gyors és egységes kialakításához \cite{tailwindUtilityClasses}. A Radix UI és a shadcn komponensek újrafelhasználható, hozzáférhetőségi szempontokat is figyelembe vevő felületi elemek készítését támogatják \cite{radixPrimitivesIntro,shadcnComponents}.
